%
% einleitung.tex -- Einleitung und Motivation
%
% (c) 2026 Baris Catan, OST Ostschweizer Fachhochschule
%
% !TEX root = ../../buch.tex
% !TEX encoding = UTF-8
%
\section{Eine harmlose Frage
\label{elliptisch:section:einleitung}}
\kopfrechts{Eine harmlose Frage}

Der Umfang einer Ellipse lässt sich nicht mit den Mitteln der
Schulmathematik berechnen.
Für den Kreis mit Radius $r$ genügt die Parametrisierung
$(x,y) = (r\cos t,\, r\sin t)$.
Wegen $\sin^2 t + \cos^2 t = 1$ fällt die Wurzel im Bogenelement weg
und es bleibt
\begin{equation}
U
=
\int_0^{2\pi} \sqrt{\dot x^2 + \dot y^2}\,dt
=
\int_0^{2\pi} r\,dt
=
2\pi r.
\label{elliptisch:equation:kreisumfang}
\end{equation}
Für die Ellipse mit den Halbachsen $a \ge b$ kann man dieselbe
Rechnung durchführen, aber die Wurzel verschwindet nicht.
Mit der numerischen Exzentrizität $k^2 = (a^2-b^2)/a^2$ bleibt
\begin{equation}
U
=
4a \int_0^{\pi/2} \sqrt{1 - k^2\sin^2 t}\,dt
\label{elliptisch:equation:ellipsenumfang}
\end{equation}
stehen \cite{elliptisch:mueller-spezielle}.
\index{Ellipse}%
\index{Exzentrizitat@Exzentrizität}%
Dieses Integral hat keine elementare Stammfunktion.
Dass das so ist, lässt sich nicht durch Probieren feststellen,
sondern nur mit algebraischen Methoden beweisen.

\subsection{Existenz und Darstellbarkeit}
Nach dem Hauptsatz der Differential- und Integralrechnung besitzt
jede stetige Funktion $f$ eine Stammfunktion, denn
$F(x) = \int_a^x f(t)\,dt$ ist differenzierbar mit $F' = f$.
Damit ist die Existenz gesichert.
Offen bleibt, ob sich $F$ als endliche Formel aus Polynomen,
Wurzeln, $\exp$, $\log$ und den trigonometrischen Funktionen
schreiben lässt.
Die Beispiele $\sqrt{ax^2+bx+c}$, $\sqrt{\tan x}$ und $e^{-x^2}$
zeigen drei mögliche Antworten:
elementar integrierbar, elementar integrierbar mit einem praktisch
unbrauchbaren Resultat, und beweisbar nicht elementar integrierbar.
\index{elementare Funktion}%
Die Frage nach der geschlossenen Darstellung ist damit keine Frage
der Analysis, sondern der Algebra.

\subsection{Der Satz von Liouville}
Der Satz von Liouville (1833/35) gibt diese algebraische Antwort.
Wenn eine elementare Funktion $f$ eine elementare Stammfunktion
besitzt, dann hat diese die Gestalt
\begin{equation}
\int f\,dx
=
v_0 + \sum_{i=1}^n c_i \log v_i,
\label{elliptisch:equation:liouville}
\end{equation}
wobei $v_0$ und die $v_i$ aus demselben Funktionenkörper stammen wie
$f$ \cite[Kapitel~8]{elliptisch:mueller-algebra}.
\index{Liouville, Satz von}%
Risch hat 1969 daraus einen Algorithmus entwickelt, der entscheidet,
ob eine elementare Stammfunktion existiert, und der heute in
Computeralgebrasystemen implementiert ist.
\index{Risch-Algorithmus}%
Für Integranden der Form $R(x,\sqrt{p(x)})$ mit rationalem $R$ und
einem Polynom $p$ hängt die Antwort vom Grad von $p$ ab:
bis Grad $2$ ist das Integral elementar, ab Grad $3$ im Allgemeinen
nicht mehr.
Das Ellipsenintegral \eqref{elliptisch:equation:ellipsenumfang}
fällt mit Grad $4$ in den zweiten Fall.
Integrale dieser Art heissen elliptische Integrale.
